\documentclass{article}

\usepackage[T1]{fontenc}

\usepackage{booktabs,graphicx,siunitx}

\begin{document}

\section{Spearman correlation}

\begin{tabular}{lllll}
\toprule
Pair & $\rho$ [95\% CI] & S & p & N \\
\midrule
satisfaction -- motivation & 0.85 [0.76, 0.90] & 1632.68 & $<$.001 & 40 \\
\bottomrule
\end{tabular}

\includegraphics{figures/01_spearman/figure_scatter.png}

$\rho$ measures how well the relationship follows a consistent up-or-down trend (monotonic), using ranks. Read sign and magnitude like Pearson's r; p tests significance. Your significance threshold ($\alpha$) is 0.05.

A Spearman correlation gave $\rho$=.85 [.76, .90], p $<$ .001, N=40.

Statistical basis: Spearman's rank correlation ($\rho$) (Spearman, C. (1904). "The proof and measurement of association between two things." American Journal of Psychology, 15(1), 72-101.).

\end{document}